%% MICE-DUO draft. Numbers in tables come from make_tables.py; \todo marks results not yet run.
\documentclass{article}

\usepackage{microtype}
\usepackage{graphicx}
\usepackage{booktabs}
\usepackage{hyperref}
\newcommand{\theHalgorithm}{\arabic{algorithm}}
\usepackage{icml2026}          % blind review
\usepackage{amsmath}
\usepackage{amssymb}
\usepackage{mathtools}
\usepackage{amsthm}
\usepackage[capitalize,noabbrev]{cleveref}
\usepackage{algorithm}
\usepackage{algorithmic}
\usepackage[textsize=tiny]{todonotes}

\theoremstyle{plain}
\newtheorem{theorem}{Theorem}
\newtheorem{proposition}{Proposition}
\newtheorem{lemma}{Lemma}
\newtheorem{corollary}{Corollary}
\theoremstyle{definition}
\newtheorem{definition}{Definition}
\newtheorem{assumption}{Assumption}
\theoremstyle{remark}
\newtheorem{remark}{Remark}

% Shared notation. Vectors bold (\mathbf), matrices calligraphic (\mathcal), names and operators upright.
\newcommand{\vx}{\mathbf{x}}
\newcommand{\vX}{\mathbf{X}}
\newcommand{\vz}{\mathbf{z}}
\newcommand{\vZ}{\mathbf{Z}}
\newcommand{\vw}{\mathbf{w}}
\newcommand{\valpha}{\boldsymbol{\alpha}}
\newcommand{\cC}{\mathcal{C}}      % context (a set of labelled rows)
\newcommand{\cY}{\mathcal{Y}}      % label set
\newcommand{\cH}{\mathcal{H}}      % hypothesis space
\newcommand{\cD}{\mathcal{D}}
\newcommand{\cS}{\mathcal{S}}      % data stream
\newcommand{\E}{\mathbb{E}}
\newcommand{\R}{\mathbb{R}}
\newcommand{\Prob}{\Pr}
\newcommand{\KL}{\mathrm{KL}}
\newcommand{\Regret}{\mathrm{Regret}}
\newcommand{\Div}{\mathrm{D}}
\DeclareMathOperator*{\argmin}{arg\,min}
\DeclareMathOperator*{\argmax}{arg\,max}
\DeclareMathOperator{\AUROC}{AUROC}

\newcommand{\mice}{\textsc{MICE}}
\newcommand{\duo}{\textsc{DUO}}
\newcommand{\method}{\textsc{MICE-DUO}}
\newcommand{\sel}{\mathrm{sel}}
\makeatletter\newcommand{\rowinput}[1]{\@@input #1 }\makeatother

\begin{document}

\twocolumn[
  \icmltitle{Select, but Safely: Relevance-Selected Contexts as Experts\\for In-Context Learning on Drifting Streams}
  \icmlsetsymbol{equal}{*}
  \begin{icmlauthorlist}
    \icmlauthor{Anonymous Author(s)}{anon}
  \end{icmlauthorlist}
  \icmlaffiliation{anon}{Anonymous Institution}
  \icmlcorrespondingauthor{Anonymous Author(s)}{anon@anon.org}
  \icmlkeywords{concept drift, tabular foundation models, in-context learning, context selection, online learning}
  \vskip 0.3in
]
\printAffiliationsAndNotice{}

\begin{abstract}
A tabular foundation model (TFM) learns a stream through its context, so the choice of which past rows to condition
on is the learning algorithm. A natural choice is \emph{relevance selection}: at every step, rank recent labelled
batches by how well they agree with the newest batch and condition on the most relevant ones, not on the most recent.
We find that relevance selection is a high-variance learner. On seven development streams it matches the accuracy of
\mice{}, a mixture of in-context experts with a worst-case safeguard, but it loses up to half a point to \mice{} on
some streams and, used alone, up to 4.6 points to a plain window. We show that this variance is not a reason to drop
selection but a reason to treat it as an expert. Inserted as one more expert into \mice{}, unchanged, the selected
context inherits \mice{}'s any-time regret bound against the plain window, and empirically it costs at most
0.15 points against \mice{} on every development stream while keeping most of its gain where concepts recur.
\todo{Mechanism result on controlled recurrence and held-out confirmation, after the pre-registered runs.}
\end{abstract}

% =====================================================================================================
\section{Introduction}
\label{sec:intro}

Tabular foundation models such as TabPFN \citep{hollmann2025tabpfn} and TabICL \citep{qu2025tabicl} predict in one
forward pass given a context of labelled rows \citep{muller2022pfn,nagler2023statistical}. On a data stream they
replace incremental learners: the model is frozen and the context is updated \citep{lourenco2026streams}. What the
model knows about the current concept is then exactly the set of rows in its context, and the stream learner reduces
to a rule for choosing that set.

Two rules are in use. A \emph{window} conditions on the most recent rows; it is safe but forgets a concept as soon as
it leaves the window. A \emph{memory} keeps past segments as stored contexts and decides which to recall; \mice{}
\citep{mice2026} mixes windows and stored segments under a two-level meta expert and guarantees, with no assumption
on the stream, that it is never much worse than the window. Both rules choose contexts made of contiguous rows.

A third rule is suggested by the analysis of non-IID batches \citep{zhao2024wnb}: the error of a learner trained on
a set of batches decomposes into a term that grows with the discrepancy between those batches and the test batch,
and a term that shrinks with their size. For an in-context learner this is the trade-off between a large context and
a relevant one. \emph{Relevance selection} resolves it per step. It ranks the recent labelled batches by the loss of
the TFM conditioned on the newest batch, keeps the most relevant ones together with a short recent anchor, and
conditions on this non-contiguous context. We call the pair of this context and the window \duo{}.

\textbf{The problem.} Relevance selection is attractive and unreliable. On the seven development streams of
\cref{sec:exp-dev} \duo{} is as accurate as \mice{} on average, gains 1.5 points on a stream whose concepts recur,
and loses about half a point on two streams where they do not; the selected context alone is 4.6 points below a
plain window on one stream (\cref{tab:dev}). Tuning the rule does not help: a pre-registered search over the size of
the recent anchor traded the losses for the gains and passed on no seed (\cref{sec:exp-anchor}).

\textbf{Our answer.} We do not choose between selection and memory. We insert the selected context as one more
expert of \mice{}, with \mice{} unchanged. The consequence is immediate on the theory side and measurable on the
empirical side.
(i)~\mice{}'s slow level is a meta expert over two predictions, the window and the fast mixture; what the fast
mixture contains does not enter its bound. The any-time discounted regret of the combined learner against the window
is therefore at most $2\ln 2$, as for \mice{} (\cref{cor:safe}).
(ii)~\mice{}'s fast level reweights its experts every batch from their last loss, so a selected context that is wrong
is down-weighted after one labelled batch under a margin, and is otherwise absorbed by the slow level.
(iii)~Against \mice{} itself there is no guarantee, and we measure it: on the development streams the combined
learner is never more than 0.15 points below \mice{} on any stream and seed and is above it on average
(\cref{tab:verdict}).

\textbf{Contributions.}
\textbf{(1)}~Relevance selection of non-contiguous contexts for TFMs on streams, derived from the discrepancy--size
trade-off of non-IID batches, and the observation that it is a high-variance learner.
\textbf{(2)}~A plug-in construction that turns any context-selection rule into an expert of \mice{} with the
window safeguard intact.
\textbf{(3)}~Pre-registered evidence: development streams, a controlled test of when selection helps, and a frozen
held-out evaluation. \todo{(3) depends on the pending runs.}

% =====================================================================================================
\section{Setting and Background}
\label{sec:setting}

\textbf{Protocol.} A stream arrives in batches $X_1,X_2,\dots$ of $B$ rows; the labels of batch $t-1$ are revealed
before batch $t$ is predicted. A frozen TFM $f_\Theta(\cdot\mid X;\cC)$ predicts class probabilities for the rows of
$X$ given a context $\cC$ of at most $M$ labelled rows. A stream learner chooses $\cC_t$ from the labelled rows of
batches $1,\dots,t-1$, or combines the predictions of several contexts.

\textbf{\mice{}.} \mice{} \citep{mice2026} keeps a set $\mathcal E$ of in-context experts: FIFO windows of 100, 300 and
$M$ rows, and up to 12 stored contexts of past segments. Its \emph{fast level} mixes all experts with weights
\begin{equation}\label{eq:fast}
  w_{t,k}\propto\exp\bigl(-\eta\,\ell^{\log}_{t-1,k}\bigr),\qquad k\in\mathcal E,
\end{equation}
where $\ell^{\log}_{t-1,k}$ is the log-loss of expert $k$ on the last labelled batch and $\eta=10$. Its \emph{slow
level} mixes the fast mixture $p^{\mathrm{fast}}_t$ with the $M$-row window $p^{\mathrm{win}}_t$ by exponential
weights on discounted half Brier scores $\ell_{t,k}\in[0,1]$,
\begin{equation}\label{eq:slow}
  L_{t,k}=\gamma L_{t-1,k}+\ell_{t,k},\quad v_{t,k}\propto e^{-\eta_2\gamma L_{t-1,k}},\quad
  p_t=v_{t,0}\,p^{\mathrm{win}}_t+v_{t,1}\,p^{\mathrm{fast}}_t,
\end{equation}
with $\gamma=0.9$ and $\eta_2=\tfrac12$.

\begin{theorem}[\citealp{mice2026}, any-time safeguard]\label{thm:outer}
For $K$ predictions combined by \eqref{eq:slow} with $\eta_2\le\tfrac12$, every sequence of labels and predictions,
every $T$ and every $k$:
$\sum_{t\le T}\gamma^{T-t}\bigl(\ell_t(p_t)-\ell_{t,k}\bigr)\le\ln K/\eta_2$.
\end{theorem}

% =====================================================================================================
\section{Relevance Selection}
\label{sec:duo}

\textbf{Motivation.} \citet{zhao2024wnb} bound the risk of a learner trained on a weighted set of non-IID batches by a
term in the discrepancy between each batch and the target distribution, weighted by the batch's share, plus a term
that decreases with the effective sample size. Small batches are individually far from the target; a large set of
them is close only if they are relevant. A TFM does not accept sample weights, and duplicating rows to emulate them
hurts \todo{cite our ResetEval negative results or drop}. We therefore implement the weights as a hard selection.

\textbf{Rule.} At step $t$ the candidates are the labelled batches $t-H,\dots,t-1$ with $H=100$. A recent anchor
$\mathcal A_t$ of the last $a$ batches is always kept. Every older candidate $b$ is scored by the mean log-loss of
the TFM conditioned on the anchor,
\begin{equation}\label{eq:score}
  s_t(b)=\tfrac1B\textstyle\sum_{(x,y)\in b}-\log f_\Theta(y\mid x;\mathcal A_t),
\end{equation}
a model-oriented discrepancy between batch $b$ and the current concept in the sense of \citet{zhao2023rdiv}. The
selected context $\cC^{\sel}_t$ is the anchor together with the $M/B-a$ lowest-scoring older batches; its prediction
is $p^{\sel}_t=f_\Theta(\cdot\mid X_t;\cC^{\sel}_t)$. All scores use labels of batches before $t$ only.

\textbf{\duo{}.} The original proposal mixes $p^{\sel}_t$ with the $M$-row window by exponential weights on
discounted log-losses ($\eta=2$, $\gamma=0.5$) and uses $a=1$ (an anchor of one batch).

\textbf{Cost.} Scoring takes one TFM call with up to $(H-a)B$ query rows, evaluated in chunks; predicting takes one
call with $B$ query rows. \todo{Measured wall-clock cost per batch relative to \mice{}.}

% =====================================================================================================
\section{\method{}: Selection as an Expert}
\label{sec:method}

\method{} is \mice{} with one more expert: $\mathcal E'=\mathcal E\cup\{\sel\}$, where expert $\sel$ predicts
$p^{\sel}_t$. Nothing else changes: the fast weights \eqref{eq:fast} now range over $\mathcal E'$, and the slow level
\eqref{eq:slow} still combines the $M$-row window with the fast mixture. No constant is re-tuned.

\begin{corollary}[The safeguard survives any added expert]\label{cor:safe}
For every stream, every selection rule and every $T$, the discounted half-Brier regret of \method{} against the
$M$-row window, and against its own fast mixture, is at most $\ln 2/\eta_2=2\ln2$.
\end{corollary}
\begin{proof}
\Cref{thm:outer} holds for arbitrary predictions $p_{t,k}$. The slow level of \method{} has $K=2$ predictions, the
window and the fast mixture over $\mathcal E'$; which experts the fast mixture contains is not used.
\end{proof}

\begin{remark}[What is not guaranteed]
\Cref{cor:safe} compares \method{} with the window, not with \mice{}. Adding an expert changes the fast mixture, and
the fast mixture of \mice{} is not one of \method{}'s slow-level experts, so a bound against \mice{} would need a third
slow-level expert. We keep \mice{} unchanged and measure the difference instead (\cref{sec:exp-dev}).
\todo{Optionally evaluate the nested variant whose slow level also contains \mice{}'s output, which has a
$\ln3/\eta_2$ bound against \mice{}; pre-register before computing.}
\end{remark}

\begin{remark}[Why a wrong selection is cheap]
Under \mice{}'s margin assumption, an expert whose loss on the last labelled batch exceeds the best by $\Delta$ has fast
weight at most $e^{-\eta\Delta}$ by \eqref{eq:fast}; with $\eta=10$, a selected context whose last log-loss exceeds
the best by 0.3 nats keeps less than 5\% of the fast weight after one labelled batch. When no expert is separated, the slow level decides whether following the fast mixture is worth
it, which is the case \cref{cor:safe} covers.
\end{remark}

\textbf{Horizon.} Relevance selection only sees the last $H$ batches; \mice{}'s stored contexts have no horizon. The
two are complementary: selection recovers concepts that recurred within $H$ batches at the granularity of single
batches, stored contexts recover older concepts at the granularity of segments.

% =====================================================================================================
\section{Experiments}
\label{sec:exp}

\textbf{Protocol.} As \mice{}: $B=100$, $M=1000$, labels one batch late, TabPFN v2 with four estimators. Every
comparison is on the same batches with the same backbone seed. All decision rules were written and committed before
the corresponding runs; we report every pre-specified variant, including failures. The 22 real streams not used in
development are held out and untouched. Backbone seeds change model randomness on fixed data and are not independent
data replicates.

\subsection{Development Streams: Selection Alone and as an Expert}
\label{sec:exp-dev}

Seven real streams (Elec2 and six streams used in \citet{mice2026}'s evaluation), first 300 batches, backbone seeds 1
and 2. \Cref{tab:dev} reports accuracy. The selected context alone is the least accurate learner on average, yet on Rialto, whose concepts recur, it is
6.3 points above FIFO and 1.2 above \mice{}. \duo{} matches \mice{} on average but is 0.50 and 0.47 points below it
on Airlines and Weather. \method{} keeps the average and removes these losses.

\begin{table}[t]
  \centering
  \caption{Accuracy (\%) on the development streams, mean over backbone seeds 1--2. sel: selected context alone;
  last column: \method{} minus \mice{}. Bold: best per stream.}
  \label{tab:dev}
  \small
  \setlength{\tabcolsep}{3.5pt}
  \begin{tabular}{lcccccr}
    \toprule
    Stream & FIFO & sel & \duo{} & \mice{} & \method{} & $\Delta$ \\
    \midrule
    \rowinput{tab_dev_rows.tex}
    \bottomrule
  \end{tabular}
\end{table}

The pre-registered criterion for \method{} asks for no negative effect: on every stream at most 0.30 points below
\mice{} and below FIFO, a mean not below \mice{}, and a mean log-loss at most 0.005 above \mice{}'s. Gains are not
required. \Cref{tab:verdict} reports the primary variant (anchor 100) and three secondary ones. The primary variant
passes on both seeds, as do the anchor-300 selection and the \duo{} mixture as experts; the anchor-500 selection fails
on seed 2 (mean $-0.05$). The criterion was written after seeing \duo{}'s per-stream results on these streams, so
this pass is exploratory; it licenses the tests of \cref{sec:exp-mech,sec:exp-heldout}, not a claim.

\begin{table*}[t]
  \centering
  \caption{Pre-registered no-negative-effect check on the development streams. Differences in accuracy points;
  $\Delta$LL: mean log-loss minus \mice{}'s.}
  \label{tab:verdict}
  \small
  \begin{tabular}{lcccccccc}
    \toprule
    Expert added to \mice{} & seed & mean vs \mice{} & wins & worst vs \mice{} & worst vs FIFO & $\Delta$LL & pass \\
    \midrule
    \rowinput{tab_verdict_rows.tex}
    \bottomrule
  \end{tabular}
\end{table*}

\subsection{Tuning Selection Alone Does Not Work}
\label{sec:exp-anchor}

Before \method{} we tried to make \duo{} safe by itself, by enlarging the recent anchor to 300 and 500 rows so that the
selection only fills the rest of the context. The rule required, on every backbone seed, a mean gain over \mice{} of
at least 0.5 points, wins on at least 5 of 7 streams and no stream more than 0.3 points below FIFO. No variant passed
on any seed. Larger anchors moved \duo{} towards FIFO: they reduced the worst loss against FIFO (from about $-0.55$ to
$-0.19$ with anchor 300) and removed the gain on Rialto (from $+1.35$ to $-0.47$ against \mice{}). Even an oracle that
picks the better of \duo{} and \mice{} per stream after the fact is only 0.23--0.28 points above \mice{}.
\todo{Move the full table to the appendix (\texttt{MiceDuo/results/night\_20261007/}).}

\subsection{When Does Selection Help? Controlled Recurrence}
\label{sec:exp-mech}

Eighteen fresh synthetic streams from \mice{}'s grid generator, unused by either method: fifteen blocks of 2000 rows,
pure real drift (shared inputs, permuted centroid labels), with $K\in\{3,5,15\}$ distinct concepts so that each
concept appears 5 times, 3 times, or once; 30 or 100 centroids; three data seeds. Pre-registered hypotheses:
(H1) no negative effect when concepts never recur; (H2) no negative effect anywhere; (H3) the gain over \mice{}
increases with recurrence ($K=3>K=5>K=15$) and is positive on at least 5 of 6 streams at $K=3$. Before the results we
noted that with blocks of 20 batches and $H=100$, a concept with $K=5$ returns exactly at the edge of the selection
horizon. \todo{Results: running.}

\subsection{Held-Out Real Streams}
\label{sec:exp-heldout}

\todo{Frozen evaluation on the 22 held-out real streams, only if \cref{sec:exp-mech} supports H1--H2; requires
explicit approval before running.}

% =====================================================================================================
\section{Related Work}
\label{sec:related}

\textbf{In-context stream learning.} TFMs with sliding memories outperform adaptive tree ensembles on drifting
streams \citep{lourenco2026streams}, and drifting priors improve them under temporal shift \citep{helli2024drift}.
\mice{} \citep{mice2026} adds a memory of stored contexts with a worst-case safeguard. Relevance selection is a third
way to build a context, at the granularity of batches and within a horizon.

\textbf{Recurring drift.} Pools of models with reuse decisions \citep{goncalves2013rcd,gama2014recurrent} and the
mixture of online and offline experts \citep{zhao2025mooe} recall trained models; drift detectors reset them
\citep{gama2004ddm}. See \citet{gama2014survey,lu2019learning} for surveys.

\textbf{Discrepancy and batch weighting.} R-divergence measures model-oriented discrepancy \citep{zhao2023rdiv};
weighting non-IID batches by discrepancy improves out-of-distribution detection \citep{zhao2024wnb}. Our score
\eqref{eq:score} is a one-batch model-oriented discrepancy used for selection rather than weighting.

\textbf{Expert advice.} Exponential weights with exp-concave losses \citep{cesabianchi2006prediction,hazan2007log}
make adding an expert cheap in regret; \cref{cor:safe} is an instance at the level of contexts.

% =====================================================================================================
\section{Limitations}

The development gain of \method{} over \mice{} is small (+0.15 points) and comes mostly from one stream. The safeguard
is against the window, not against \mice{}. Selection costs extra TFM calls per batch. Backbone seeds share data.
\todo{Update after mechanism and held-out results.}

\section{Conclusion}

Relevance selection gives an in-context learner the right past rows when concepts recur and the wrong ones when they
do not. As an expert of a mixture with a window safeguard, the first property is kept and the second is contained.
\todo{Final claim after pending results.}

\bibliography{refs}
\bibliographystyle{icml2026}

\end{document}
